\documentclass[12pt]{article}

\usepackage{amsmath, amssymb}
\usepackage{geometry}
\usepackage{booktabs}
\usepackage{array}
\usepackage{hyperref}
\usepackage{parskip}
\usepackage{microtype}

\geometry{margin=1.2in}

\title{AI Reconstructions of Newton's Gravitational Reasoning:\\
From Kepler's Laws to Universal Gravitation}
\author{}
\date{}

\begin{document}

\maketitle

\begin{abstract}
Several recent research programs have asked whether a modern AI system,
given observational data of the kind available to Newton, could independently
rediscover the law of universal gravitation.
This note reviews those efforts, situates them within Newton's original
mathematical reasoning in the \textit{Principia}, and draws comparisons
between Newton's geometric method, modern calculus, symbolic regression,
and Bayesian inference.
The central finding is that Newton's conceptual chain---area law implies
central force; harmonic law implies inverse-square dependence; inverse-square
force implies conic orbits---can be reproduced algorithmically, though the
epistemological routes differ substantially.
\end{abstract}

%----------------------------------------------------------------------
\section{Introduction}

Newton's derivation of universal gravitation is often summarized as
``fitting an inverse-square law to Kepler's laws,'' but the actual structure
of the argument in \textit{Principia Mathematica} (1687) is considerably
richer.
Newton did not curve-fit data; he built a geometric dynamical theory and
demonstrated that Kepler's three laws follow as logical consequences of a
single universal principle.

Two broad categories of modern work revisit this achievement computationally:

\begin{enumerate}
  \item \textbf{Deep-learning systems} that infer Newtonian gravity from raw
        trajectory data.
  \item \textbf{Symbolic regression systems} that search for closed-form force
        laws consistent with observed dynamics.
\end{enumerate}

Together they provide a modern lens on a classic question: how much of
Newton's achievement was data-driven, and how much was conceptual innovation?

%----------------------------------------------------------------------
\section{Newton's Derivation in the \textit{Principia}}

\subsection{Step A: Kepler's Area Law Implies a Central Force}

Kepler's second law states that a planet sweeps equal areas in equal times.
Newton reinterpreted this kinematically.
If no force acts, a body moves in a straight line (Axiom~I).
Newton showed (Proposition~1, Book~I) by a limiting triangular-area argument
that, if impulses always point toward a fixed center, the resulting polygonal
path sweeps equal triangular areas in equal times.
In the limit of continuous force, equal areas in equal times is therefore
\emph{equivalent} to the force being directed toward a fixed point.

\textbf{Conclusion:} Kepler's second law implies a \emph{central force}.

\subsection{Step B: Kepler's Harmonic Law Implies an Inverse-Square Force}

Kepler's third law states
\[
  T^{2} \propto a^{3},
\]
where $T$ is the orbital period and $a$ is the semi-major axis.
For a circular orbit of radius $r$, the centripetal acceleration is
\[
  a_{r} = \frac{v^{2}}{r}, \qquad v = \frac{2\pi r}{T}.
\]
Substituting and applying Kepler's third law ($T^{2} = k\,r^{3}$),
\[
  a_{r} = \frac{4\pi^{2} r}{T^{2}} = \frac{4\pi^{2} r}{k\,r^{3}}
         = \frac{4\pi^{2}}{k}\,\frac{1}{r^{2}}.
\]
Hence the centripetal acceleration---and therefore the gravitational
force---falls off as $r^{-2}$.

\textbf{Conclusion:} Kepler's third law implies an \emph{inverse-square
central force}.

\subsection{Step C: An Inverse-Square Force Produces Elliptical Orbits}

Newton then reversed the logic (Proposition~11, Book~I): assuming a central
force proportional to $1/r^{2}$, he proved geometrically that the orbit is a
conic section---an ellipse when the total energy is negative, a parabola when
it is zero, and a hyperbola when positive.
This established that the inverse-square law is \emph{equivalent} to Kepler's
three laws, not merely consistent with them.

\subsection{Step D: Universality---The Moon Test}

Newton compared the gravitational acceleration at Earth's surface with the
centripetal acceleration of the Moon inferred from its orbital motion.
Let $r_{\oplus}$ and $r_{\text{Moon}}$ denote the radii of the Earth and the
Moon's orbit respectively.
If gravity is inverse-square,
\[
  \frac{g_{\text{surface}}}{a_{\text{Moon}}}
  = \left(\frac{r_{\text{Moon}}}{r_{\oplus}}\right)^{2}.
\]
Newton's first calculation disagreed because he used an inaccurate value for
Earth's radius.
When Jean Picard's new measurement (1671) became available in England,
Newton recomputed the ratio and found agreement to within observational error.
This was the conceptual leap from \emph{solar} gravity to
\textbf{universal gravitation}: the same force that accelerates a falling
body near Earth's surface governs the Moon's orbit.

%----------------------------------------------------------------------
\section{Modern Calculus Reconstruction}

The following is a calculus-based restatement of Newton's reasoning,
using notation compatible with the AI reconstructions discussed in
Section~\ref{sec:ai}.

\subsection{Area Law and Angular Momentum Conservation}

Kepler's second law in modern notation:
\[
  \frac{dA}{dt} = \tfrac{1}{2}r^{2}\dot\theta = \text{constant} \equiv \tfrac{h}{2},
\]
where $h = r^{2}\dot\theta$ is the \emph{specific angular momentum}
(angular momentum per unit mass).
Equivalently,
\[
  L = m\,r^{2}\dot\theta = m\,h = \text{constant}.
\]
The constancy of $h$ means the net tangential acceleration vanishes, i.e.\
the force is central.

\subsection{Harmonic Law and the Inverse-Square Force}

For a circular orbit the centripetal acceleration satisfies
$a_{r} = v^{2}/r = 4\pi^{2}r/T^{2}$.
Kepler's third law gives $T^{2} = k\,r^{3}$, so
\[
  a_{r} = \frac{4\pi^{2}}{k}\,\frac{1}{r^{2}}, \qquad
  F = m\,a_{r} \propto \frac{1}{r^{2}}.
\]

\subsection{Orbit Equation for a $1/r^{2}$ Force}

Substituting $u = 1/r$ and using $h = r^{2}\dot\theta$, the radial equation
of motion becomes the \emph{Binet equation},
\[
  \frac{d^{2}u}{d\theta^{2}} + u = \frac{GM}{h^{2}},
\]
where $M$ is the mass of the attracting body and $G$ is the gravitational
constant.
The general solution is
\[
  u(\theta) = \frac{GM}{h^{2}}\bigl[1 + e\cos(\theta - \theta_{0})\bigr],
\]
or equivalently,
\[
  r(\theta) = \frac{p}{1 + e\cos\theta}, \qquad
  p = \frac{h^{2}}{GM},
\]
which is the polar equation of a conic section with semi-latus rectum $p$
and eccentricity $e$.
Kepler's third law follows directly: integrating the area swept over one
orbit,
\[
  T^{2} = \frac{4\pi^{2}}{GM}\,a^{3}.
\]
This is Kepler's third law with an explicit physical constant $GM$;
the empirical proportionality constant $k$ is identified as $4\pi^{2}/(GM)$.

\subsection{Comparison: Newton's Geometry vs.\ Modern Calculus}

\begin{table}[h]
\centering
\renewcommand{\arraystretch}{1.3}
\begin{tabular}{p{3cm} p{5.5cm} p{5.5cm}}
\toprule
\textbf{Concept} & \textbf{Newton's Geometric Method} & \textbf{Modern Calculus} \\
\midrule
Area law
  & Equal triangular areas via limiting polygon argument
  & $h = r^{2}\dot\theta = \text{const}$ \\
Central force
  & Force always directed toward the focus
  & No tangential acceleration: $\ddot\theta + 2\dot r\dot\theta/r = 0$ \\
Harmonic law
  & Period--radius ratio for circular orbits
  & $T^{2} = (4\pi^{2}/GM)\,a^{3}$ \\
Orbit shape
  & Conic sections from geometric lemmas (Prop.\ 11, Book~I)
  & Solution to the Binet equation \\
Inverse-square law
  & Derived from harmonic law for circular orbits
  & Identified from the Binet equation \\
\bottomrule
\end{tabular}
\caption{Newton's geometric method and its modern calculus counterpart.}
\end{table}

Newton wrote in synthetic geometry largely because calculus was not yet
widely accepted and he wanted the logical certainty of Euclidean proof.
Modern reconstructions use differential equations because they are more
compact, extend naturally to perturbations, and align with how AI systems
represent dynamics.

%----------------------------------------------------------------------
\section{AI Reconstructions of Newtonian Gravity}
\label{sec:ai}

\subsection{Deep Learning: Lemos et al.}

Lemos, Jeffrey, Cranmer, Battaglia, and Ho
trained a neural network whose architecture encodes the inductive bias of
classical mechanics (pairwise interactions, conservation laws) and fit it to
roughly 30 years of real Solar System trajectory data.
The network learned:
\begin{itemize}
  \item the functional form of the gravitational force law,
  \item the relative masses of planets (which are not directly observed).
\end{itemize}
Symbolic regression applied to the learned force representation recovered
$F \propto 1/r^{2}$, matching Newton's law.
The network's acceleration vectors pointed toward the Sun (central force),
and the extracted force law was consistent with the full inverse-square form.
This is one of the clearest demonstrations that a modern AI system can
reconstruct Newton's dynamical model from observational trajectory data
alone, without prior physical knowledge.

\subsection{Concept-Driven Discovery: AI-Newton}

Fang, Jian, Li, and Ma (Peking University) built a system called
\textit{AI-Newton} that takes raw, multi-experiment data, proposes
interpretable physical concepts, and generalizes them into universal laws.
Starting from trajectory data and without being given Kepler's laws or
Newton's laws, the system:
\begin{itemize}
  \item identified central-force behavior from orbital data,
  \item generalized the force law into a symbolic form consistent with
        inverse-square gravity,
  \item rediscovered Newton's second law, conservation of energy, and
        universal gravitation.
\end{itemize}
The conceptual progression---empirical regularities to central force to
inverse-square law---mirrors Newton's own chain of reasoning, executed
without human supervision.

\subsection{Symbolic Regression: Eureqa and AI-Feynman}

Symbolic regression systems such as Eureqa (Schmidt \& Lipson, 2009) and
AI-Feynman (Udrescu \& Tegmark, 2020) search the space of mathematical
expressions for formulas that minimize prediction error while penalizing
complexity.
Eureqa has been used to recover Newtonian force laws from simulated
dynamical data; AI-Feynman has systematically rediscovered equations from
the Feynman Lectures on Physics.
These systems use the structure of the search space (powers, products,
inverses) analogously to Newton's geometric constructions, and both converge
on the inverse-square law from trajectory data.

\subsection{Comparison: Newton's Reasoning vs.\ AI Reconstruction}

\begin{table}[h]
\centering
\renewcommand{\arraystretch}{1.3}
\begin{tabular}{p{3.5cm} p{4.5cm} p{4.5cm}}
\toprule
\textbf{Step} & \textbf{Newton} & \textbf{AI/ML Analogue} \\
\midrule
Elliptical orbits
  & Orbits arise from central force plus inertia
  & Neural networks learn orbital trajectories from data \\
Equal areas
  & Implies central force (Prop.\ 1, Book~I)
  & Learned acceleration vectors point toward the Sun \\
$T^{2} \propto a^{3}$
  & Implies inverse-square force
  & Symbolic regression recovers $F \propto 1/r^{2}$ \\
Universality
  & Moon test connects orbital and terrestrial gravity
  & Multi-body networks infer $GM$ and mass ratios simultaneously \\
\bottomrule
\end{tabular}
\caption{Newton's reasoning and its AI analogues.}
\end{table}

%----------------------------------------------------------------------
\section{Perturbations and Multi-Body Dynamics}

Newton recognized that Kepler's laws are approximate because planets
mutually perturb one another, and he handled this with geometric
approximations and early perturbation arguments.
Modern AI systems address the same problem differently: a network trained on
Solar System data observes that Mars's orbit deviates from a perfect ellipse
and that these deviations correlate with Jupiter's position.
The network learns a vector field of accelerations that includes all bodies
simultaneously, without perturbation theory or series expansions.
Symbolic regression applied to the learned dynamics can then recover the
full $N$-body law,
\[
  \ddot{\mathbf{r}}_{i} = \sum_{j \neq i} \frac{G m_{j}}{r_{ij}^{2}}\,\hat{\mathbf{r}}_{ij},
\]
which Newton could only approximate analytically.

%----------------------------------------------------------------------
\section{Newton's Approach vs.\ Bayesian Inference}

Newton's reasoning was model-first: derive a theoretical law from geometric
principles, then check whether the data are consistent with it, attributing
discrepancies to measurement error or perturbations.
This is the opposite of modern curve-fitting; Newton trusted the mathematics
more than the data.

A Bayesian physicist would instead assign a prior probability distribution
over force laws, compute likelihoods from the observational data, and
update beliefs via Bayes' rule.
If planetary trajectory data are fed into a Bayesian model-selection
framework, the inverse-square law emerges as the overwhelmingly preferred
model.

The observation that Newton's geometric reasoning, Bayesian statistical
inference, and AI symbolic regression all converge on $F \propto 1/r^{2}$
through entirely different epistemological routes is striking: it suggests
that the inverse-square law is, in a precise sense, the simplest explanation
of the available data.

%----------------------------------------------------------------------
\section{Conclusions}

These projects demonstrate that:
\begin{itemize}
  \item Newton's derivation of universal gravitation is algorithmically
        reproducible from trajectory data.
  \item Kepler's empirical laws contain sufficient information for an AI to
        infer the same universal law Newton did.
  \item The conceptual chain---area law, harmonic law, inverse-square
        force, conic orbits---can be recovered without prior physics
        knowledge, whether by symbolic regression, deep learning, or
        Bayesian model selection.
\end{itemize}

At the same time, important differences remain.
Newton unified terrestrial and celestial gravity through a physical argument
(the Moon test) that required conceptual ingenuity as well as data.
Current AI systems recover the force law from trajectory data but require
human judgment to interpret the result as ``universal'' gravitation rather
than a local fit.
The question of how much of Newton's achievement was data-driven and how
much was conceptual innovation remains partially open---and these AI
studies suggest that both components can be partially automated.

%----------------------------------------------------------------------
\begin{thebibliography}{9}

\bibitem{lemos2023}
P.~Lemos, N.~Jeffrey, M.~Cranmer, P.~Battaglia, and S.~Ho,
``Rediscovering Newton's Gravity and Solar System Properties Using Deep
Learning and Inductive Biases,''
\textit{Machine Learning: Science and Technology}, 2023.

\bibitem{ainewton}
Y.-L.~Fang, D.-S.~Jian, X.~Li, and Y.-Q.~Ma,
``AI-Newton: A Concept-Driven Physical Law Discovery System Without Prior
Physical Knowledge,''
Peking University preprint, 2023.

\bibitem{eureqa}
M.~Schmidt and H.~Lipson,
``Distilling Free-Form Natural Laws from Experimental Data,''
\textit{Science}, vol.~324, pp.~81--85, 2009.

\bibitem{aifeynman}
S.-M.~Udrescu and M.~Tegmark,
``AI Feynman: A Physics-Inspired Method for Symbolic Regression,''
\textit{Science Advances}, vol.~6, no.~16, 2020.

\bibitem{principia}
I.~Newton, \textit{Philosophi\ae{} Naturalis Principia Mathematica},
London, 1687.

\end{thebibliography}

\end{document}
